\subsection{Setup and Formal Definition of the Tokenisation Operator}

To rigorously understand the failure modes of patch-based architectures like Chronos-Bolt,
we must first formally define the tokenisation process.
Let $x[n]$ represent a discrete-time, univariate signal where the index $n \in \mathbb{Z}$
is sampled at a continuous frequency of $f_s$.

Unlike classical autoregressive models that process temporal data one scalar time-step at a
time, modern foundation models employ \emph{patch-based tokenisation}.
This technique groups contiguous samples to radically reduce the effective sequence length
and improve the computational efficiency of the transformer's attention mechanisms.
This is achieved by sliding a window of fixed size $P$ (the \emph{patch length}) across the
signal, moving forward by a step size of $S$ (the \emph{stride}).

We can extract this sequence of tokens as follows:
\begin{equation}
	\mathbf{p}_k
	= \bigl(x[kS],\; x[kS+1],\; \ldots,\; x[kS+P-1]\bigr)
	\in \mathbb{R}^P,
	\qquad k \in \mathbb{N}.
\end{equation}

\begin{example}[Constructing patches by hand]
	Let $x = [x_0, x_1, x_2, x_3, x_4, x_5, x_6, x_7]$ with $P = 4$ and $S = 2$.
	The first three patches are:
	\[
	\mathbf{p}_0 = [x_0,x_1,x_2,x_3], \qquad
	\mathbf{p}_1 = [x_2,x_3,x_4,x_5], \qquad
	\mathbf{p}_2 = [x_4,x_5,x_6,x_7].
	\]
	Because $S < P$, consecutive patches overlap by $P - S = 2$ samples
	(e.g.\ $\mathbf{p}_0$ and $\mathbf{p}_1$ share $x_2$ and $x_3$).
	With $S = P = 4$ (no overlap), the patches would instead be
	$\mathbf{p}_0{=}[x_0,\ldots,x_3]$, $\mathbf{p}_1{=}[x_4,\ldots,x_7]$,
	with no shared samples at all.
\end{example}

\begin{definition}[Patch Tokenisation Operator]
	We can formalise this extraction as a linear map acting on the bounded signal space.
	The map $\mathcal{T}: \ell^\infty(\mathbb{Z}) \to (\mathbb{R}^P)^{\mathbb{N}}$,
	as defined above, is referred to as the \emph{patch tokenisation operator}.
	When collecting an arbitrary number of $K$ patches from a signal vector $\mathbf{x}$,
	the operation can be represented as a matrix multiplication:
	\begin{equation}
		\mathbf{P} = \mathcal{T}\{x\} = \mathbf{W} \cdot \mathbf{x}.
	\end{equation}
	Here, $\mathbf{W} \in \{0,1\}^{KP \times N}$ is a \emph{block-binary selection matrix}.
	Its $k$-th block consists of $P$ rows that act as a mask, picking out exactly $P$
	consecutive signal samples starting at the temporal index $kS$.
\end{definition}

\begin{example}[Explicit matrix form]
	Take $N = 6$, $P = 3$, $S = 3$ (non-overlapping), giving $K = 2$ patches.
	The selection matrix and the resulting patch matrix are:
	\[
	\mathbf{W}
	= \begin{pmatrix}
		1 & 0 & 0 & 0 & 0 & 0 \\
		0 & 1 & 0 & 0 & 0 & 0 \\
		0 & 0 & 1 & 0 & 0 & 0 \\
		0 & 0 & 0 & 1 & 0 & 0 \\
		0 & 0 & 0 & 0 & 1 & 0 \\
		0 & 0 & 0 & 0 & 0 & 1 \\
	\end{pmatrix},
	\qquad
	\mathbf{P}
	= \mathbf{W}\mathbf{x}
	= \begin{pmatrix} x_0 \\ x_1 \\ x_2 \\ x_3 \\ x_4 \\ x_5 \end{pmatrix},
	\]
	which, when reshaped into $K \times P = 2 \times 3$ block form, gives
	$\mathbf{p}_0 = [x_0,x_1,x_2]$ and $\mathbf{p}_1 = [x_3,x_4,x_5]$.
	Each row of $\mathbf{W}$ contains exactly one non-zero entry (a 1) at the position of
	the corresponding signal sample, confirming the pure selection nature of the operator.
\end{example}

\noindent The operator $\mathcal{T}$ is \emph{linear} because:
\begin{enumerate}
	\item $\mathcal{T}\{ax\} = a\,\mathcal{T}\{x\}$ for any scalar $a$ — scaling the
	signal scales every patch proportionally.
	\item $\mathcal{T}\{x+y\} = \mathcal{T}\{x\} + \mathcal{T}\{y\}$ — the patch of a
	sum equals the sum of the patches.
\end{enumerate}
Linearity is the key property that makes the rank-collapse analysis in
Section~\ref{sec:one} tractable.

% -----------------------------------------------------------------
\subsection{The Mechanics of Phase-Locking and Degeneracy}

The fundamental vulnerability of the patch tokenisation operator emerges when the periodic
structure of the input signal aligns perfectly with the model's fixed stride grid.
We refer to this specific synchronisation as \emph{phase-locking}.

\begin{proposition}[Condition for Token Identity]
	Two consecutive patches, $\mathbf{p}_k$ and $\mathbf{p}_{k+1}$, will be
	mathematically identical if and only if the signal amplitude remains constant across
	the stride shift for every element within the observation window:
	\begin{equation}
		x[n] = x[n + S] \quad \forall\, n \in \{kS,\, kS+1,\, \ldots,\, kS+P-1\}.
	\end{equation}
	Globally, if the underlying continuous signal is a pure sinusoid (or any periodic
	signal) with a fundamental period of $T_0$ samples, this exact alignment occurs
	whenever the stride $S$ spans an integer multiple of the signal's period:
	\begin{equation}
		\boxed{S = m \cdot T_0, \quad m \in \mathbb{Z}^+.}
	\end{equation}
	Substituting $T_0 = f_s / f$, we can express this condition directly in terms of the
	signal's continuous frequency $f$ and the sampling frequency $f_s$:
	\begin{equation}
		f = m \cdot \frac{f_s}{S}, \qquad m \in \mathbb{Z}^+.
	\end{equation}
\end{proposition}

\begin{proof}
	To satisfy $\mathbf{p}_k = \mathbf{p}_{k+1}$ strictly component-wise, it is required
	that $x[kS + j] = x[(k+1)S + j]$ for every internal patch index
	$j \in \{0,\ldots,P-1\}$.
	For a periodic signal with $T_0 = f_s/f$, this strict equality holds whenever the
	stride shift $S$ represents a traversal of exactly $m$ full wave cycles.
	In such instances, the sliding window captures the geometrically identical segment of
	the waveform at every step, and the proposition follows.
\end{proof}

\begin{example}[Phase-locking with a discrete sinusoid]
	Let $f_s = 8\,\text{Hz}$, $f = 2\,\text{Hz}$, so $T_0 = f_s/f = 4$ samples.
	Choose $S = 4$ ($= 1 \cdot T_0$) and $P = 4$.  The signal values are:
	\begin{center}
		\begin{tabular}{c|cccccccc}
			$n$    & 0 & 1 & 2 & 3 & 4 & 5 & 6 & 7 \\\hline
			$x[n]$ & $0$ & $1$ & $0$ & $-1$ & $0$ & $1$ & $0$ & $-1$
		\end{tabular}
	\end{center}
	Then $\mathbf{p}_0 = [0,1,0,-1]$ and $\mathbf{p}_1 = [0,1,0,-1]$:
	the two patches are \textbf{identical}.
	If instead we set $S = 8 = 2T_0$, every patch still starts at phase $0$, and the
	result is the same.
	Conversely, with $S = 3$ (not a multiple of $T_0 = 4$), we obtain
	$\mathbf{p}_0 = [0,1,0,-1]$ and $\mathbf{p}_1 = [-1,0,1,0]$, which are distinct —
	no phase-locking occurs.
\end{example}

% -----------------------------------------------------------------
\subsection{The Consequence: Loss of Observability}

When the phase-locking condition $f = m \cdot (f_s/S)$ is met, we achieve a state where
$\mathbf{p}_k = \mathbf{p}_{k+1}$ for all $k$.
The token matrix $\mathbf{P}$ undergoes a mathematical \emph{rank collapse}, reducing its
rank to exactly 1.

This creates a \emph{degenerate token sequence}, which induces severe representational
failures within the network:

\begin{itemize}
	\item \textbf{Complete Temporal Blindness.}
	Because the transformer processes a sequence of completely identical vectors, the
	temporal variation of the actual underlying signal becomes entirely invisible to
	the model's self-attention mechanisms.
	The model cannot perceive that a wave is oscillating; it simply perceives a
	constant, flat embedding.
	
	\item \textbf{Observational Equivalence and Aliasing.}
	Any two distinct signals that differ only by a phase shift of $S$ samples produce
	the exact same sequence of tokens.
	They become observationally indistinguishable to the decoder, making it
	impossible for the model to correctly reconstruct or forecast the true phase of
	the original signal.
	
	\smallskip
	\begin{example}[Two signals, one token sequence]
		With $f_s = 8\,\text{Hz}$, $f = 2\,\text{Hz}$, $S = 4$, consider two
		phase-shifted signals:
		\[
		x_1[n] = \cos\!\left(\tfrac{2\pi \cdot 2\,n}{8}\right),
		\qquad
		x_2[n] = \cos\!\left(\tfrac{2\pi \cdot 2\,n}{8} + \tfrac{\pi}{2}\right).
		\]
		$x_1$ has values $[1, 0, -1, 0, 1, 0, -1, 0, \ldots]$ and $x_2$ has values
		$[0, -1, 0, 1, 0, -1, 0, 1, \ldots]$.
		Their patch sequences with $P = 4$:
		\begin{align*}
			x_1 &: \mathbf{p}_0{=}[1,0,-1,0],\quad
			\mathbf{p}_1{=}[1,0,-1,0],\quad \ldots \\
			x_2 &: \mathbf{p}_0{=}[0,-1,0,1],\quad
			\mathbf{p}_1{=}[0,-1,0,1],\quad \ldots
		\end{align*}
		Each signal individually produces a \emph{constant} token sequence.
		Although the two constant sequences differ from each other here, the decoder
		cannot infer the phase trajectory over time from either one —
		both look as if the signal is not evolving.
	\end{example}
	
	\item \textbf{Soft Degradation at Boundary Frequencies.}
	This failure is not isolated to exact integer multiples.
	Even when patches are not perfectly identical but merely highly correlated
	(i.e.\ when $S \approx m T_0$), the sequence suffers a \emph{soft} loss of
	observability.
	The effective information bandwidth is severely restricted, explaining why
	predictive accuracy degrades smoothly around these critical frequency boundaries
	before fully collapsing at the exact harmonic frequencies.
\end{itemize}

\noindent For example, in the specific architecture of Chronos-Bolt, the sampling frequency
is $f_s = 512\,\text{Hz}$ and the stride is $S = 16$.
According to our derived condition $f = m \cdot (512/16)$, the fundamental patch frequency
is precisely $32\,\text{Hz}$.
Consequently, any input signal at $32\,\text{Hz}$, $64\,\text{Hz}$, $96\,\text{Hz}$, or
other integer harmonics will trigger this structural aliasing, resulting in an immediate
and sharp collapse in the model's predictive accuracy.
